% ============================================================
%  Diffusion Policy for Robot Manipulation — Conference Paper
%  ML 6140, Northeastern University, April 2026
%  Format: NeurIPS / ICML style (10 pt, single column)
% ============================================================
\documentclass[10pt,letterpaper]{article}

% ── Core packages ────────────────────────────────────────────
\usepackage[T1]{fontenc}
\usepackage[utf8]{inputenc}
\usepackage{microtype}          % Improved typography

% ── Fonts: Times-like serif body + matching math ─────────────
\usepackage{newtxtext}
\usepackage{newtxmath}

% ── Page geometry (NeurIPS-style margins) ────────────────────
\usepackage[
  letterpaper,
  top=1in, bottom=1.125in,
  left=1.25in, right=1.25in,
  headheight=14pt
]{geometry}

% ── Math ─────────────────────────────────────────────────────
\usepackage{amsmath}
\usepackage{bm}                 % Bold math symbols

% ── Graphics & figures ───────────────────────────────────────
\usepackage{graphicx}
\usepackage{subcaption}
\usepackage{wrapfig}
\graphicspath{{./}}

% ── Tables ───────────────────────────────────────────────────
\usepackage{booktabs}
\usepackage{array}
\usepackage{multirow}
\usepackage{tabularx}
\usepackage{xcolor}             % For table row colours
\usepackage{colortbl}

% Palette
\definecolor{HeaderBlue}{RGB}{30,58,95}
\definecolor{RowAlt}{RGB}{240,245,252}
\definecolor{BestGreen}{RGB}{198,239,206}
\definecolor{CodeRed}{RGB}{178,34,34}
\definecolor{AccentBlue}{RGB}{74,144,217}

% ── Code listings ────────────────────────────────────────────
\usepackage{listings}
\lstset{
  basicstyle=\ttfamily\small,
  breaklines=true,
  frame=single,
  rulecolor=\color{gray!40},
  backgroundcolor=\color{gray!6},
  keywordstyle=\color{blue!70}\bfseries,
  commentstyle=\color{green!40!black},
  stringstyle=\color{CodeRed},
  showstringspaces=false,
  columns=fullflexible,
  keepspaces=true,
}

% ── Lists ────────────────────────────────────────────────────
\usepackage{enumitem}
\setlist[itemize]{leftmargin=1.5em, itemsep=1pt, topsep=3pt}
\setlist[enumerate]{leftmargin=1.5em, itemsep=1pt, topsep=3pt}

% ── Section heading style ────────────────────────────────────
\usepackage{titlesec}
\titleformat{\section}{\large\bfseries\color{HeaderBlue}}{{\thesection.}}{0.5em}{}
\titleformat{\subsection}{\normalsize\bfseries}{{\thesubsection}}{0.5em}{}
\titleformat{\subsubsection}{\normalsize\itshape\bfseries}{{\thesubsubsection}}{0.5em}{}
\titlespacing*{\section}{0pt}{14pt plus 2pt minus 2pt}{6pt plus 1pt}
\titlespacing*{\subsection}{0pt}{10pt plus 2pt minus 2pt}{4pt plus 1pt}
\titlespacing*{\subsubsection}{0pt}{8pt plus 1pt minus 1pt}{3pt plus 1pt}

% ── Captions ─────────────────────────────────────────────────
\usepackage[
  font=small, labelfont=bf,
  labelsep=period,
  justification=justified
]{caption}

% ── Headers / footers ────────────────────────────────────────
\usepackage{fancyhdr}
\pagestyle{fancy}
\fancyhf{}
\renewcommand{\headrulewidth}{0.4pt}
\fancyhead[L]{\small\itshape Diffusion Policy for Robot Manipulation}
\fancyhead[R]{\small\itshape ML 6140 --- Northeastern University}
\fancyfoot[C]{\small\thepage}

% ── Bibliography ─────────────────────────────────────────────
\usepackage{natbib}
\bibpunct{[}{]}{,}{n}{}{;}

% ── Hyperlinks ───────────────────────────────────────────────
\usepackage{hyperref}
\hypersetup{
  colorlinks=true,
  linkcolor=HeaderBlue,
  citecolor=HeaderBlue,
  urlcolor=AccentBlue,
  pdftitle={Diffusion Policy for Robot Manipulation},
  pdfauthor={Gupta, Sakhamuru, Maligireddy},
}

% ── Utility macros ───────────────────────────────────────────
\newcommand{\Tobs}{T_{\text{obs}}}
\newcommand{\Tpred}{T_{\text{pred}}}
\newcommand{\Tact}{T_{\text{act}}}
\newcommand{\alphabar}{\bar{\alpha}}
\newcommand{\eps}{\boldsymbol{\varepsilon}}
\newcommand{\dataset}{\mathcal{D}}
\newcommand{\method}[1]{\textbf{#1}}

% ── Paragraph spacing ────────────────────────────────────────
\setlength{\parskip}{4pt}
\setlength{\parindent}{10pt}

% ============================================================
\begin{document}

% ── TITLE ────────────────────────────────────────────────────
\begin{center}
  {\LARGE\bfseries\color{HeaderBlue}
    Diffusion Policy for Robot Manipulation: \\[4pt]
    Learning to Push with DDPM, DDIM, and Flow Matching
  }\par\vspace{12pt}
  {\large
    Yashvardhan Gupta\textsuperscript{$\dagger$}\quad
    Vineeth Sakhamuru\textsuperscript{$\dagger$}\quad
    Sai Krishna Reddy Maligireddy\textsuperscript{$\dagger$}
  }\par\vspace{6pt}
  {\normalsize
    \textsuperscript{$\dagger$}Northeastern University, Khoury College of Computer Sciences\\
    \texttt{\{gupta.yashv, sakhamuru.v, maligireddy.s\}@northeastern.edu}
  }\par\vspace{6pt}
  {\small\itshape ML 6140: Machine Learning \quad|\quad April 2026}
\end{center}

\vspace{8pt}
\noindent\rule{\linewidth}{1.2pt}
\vspace{-4pt}

% ── ABSTRACT ─────────────────────────────────────────────────
\begin{center}
  \begin{minipage}{0.92\linewidth}
    \textbf{Abstract.}\;
    We implement and evaluate \textit{Diffusion Policy}~\citep{chi2023diffusion} for
    robot manipulation on the PushT benchmark---a contact-rich 2D task requiring a
    circular end-effector to push a T-shaped block onto a target region.
    The central hypothesis is that treating action generation as an iterative
    \textit{denoising} problem enables policies to capture the multi-modal structure
    of expert demonstrations, overcoming the mode-averaging failure of standard
    Behavioral Cloning~(BC).
    We train a 1D temporal U-Net (68.95M parameters) with Feature-wise Linear
    Modulation~(FiLM) conditioning under three formulations: DDPM (stochastic, 100
    steps), DDIM (deterministic, 10 steps), and Flow Matching (ODE-based, 10 steps).
    All three diffusion methods substantially outperform the BC baseline (4\%
    success), achieving 80\%, 92\%, and \textbf{98\%} success rates over 50 evaluation
    episodes---matching or exceeding the published benchmarks of~\citet{chi2023diffusion}.
    We additionally contribute a DDIM numerical-stability fix that lifts success from
    0\% to 92\%, and a DDIM inference-step ablation demonstrating that
    \textbf{5 denoising steps achieve 100\% success}, revealing a
    surprisingly well-conditioned denoising landscape.
    Our results confirm that diffusion's advantage over BC stems primarily from
    temporal coherence in action prediction rather than multi-modality alone.
  \end{minipage}
\end{center}

\vspace{-2pt}
\noindent\rule{\linewidth}{1.2pt}
\vspace{4pt}

% ============================================================
\section{Introduction}
% ============================================================

Robot learning from demonstration faces a fundamental challenge: expert behavior
is inherently \textit{multi-modal}.
From the same initial configuration, a skilled human operator might approach a
target object from the left or from the right---both strategies are correct, yet
they correspond to completely different action trajectories.
Standard behavioral cloning~(BC) trains a deterministic regressor mapping
observations to actions.
When multiple valid actions exist, the regressor predicts their mean---an action
that is \emph{wrong in every mode}~\citep{pomerleau1989alvinn}.
This failure, known as \textit{mode averaging}, fundamentally limits BC in
multi-modal environments.

Diffusion models~\citep{ho2020ddpm} offer a principled alternative.
Rather than regressing to a single action, they learn a probability
\textit{distribution} over actions and sample from it via an iterative denoising
process.
At inference time, starting from Gaussian noise and running the learned reverse
process produces a sample from the data distribution, capturing all modes rather
than averaging them.
\citet{chi2023diffusion} demonstrated that this approach---\textit{Diffusion Policy}---achieves
state-of-the-art results on a range of robot manipulation tasks, including the
PushT benchmark we study here.

In this work, we implement Diffusion Policy from scratch and conduct a systematic
evaluation of three inference methods on PushT:
(1)~\method{DDPM}---the original stochastic reverse process with 100 denoising steps;
(2)~\method{DDIM}---a deterministic, skip-step variant requiring only 10 network
evaluations~\citep{song2021ddim};
and (3)~\method{Flow Matching}---an alternative continuous-time formulation that
replaces the noisy diffusion trajectory with straight-line interpolation between
data and noise~\citep{lipman2023flow,liu2023rectified}.

\paragraph{Contributions.}
\begin{itemize}[itemsep=3pt]
  \item \textbf{Full from-scratch implementation:} ConditionalUnet1D with FiLM
    conditioning, all three samplers, EMA training, receding-horizon evaluation
    harness, and 120 unit tests.
  \item \textbf{Critical DDIM numerical fix:} identification and correction of a
    cosine-schedule denominator blow-up that caused 0\% success despite a
    well-trained model; the fix recovers 92\% success.
  \item \textbf{Comprehensive evaluation:} 50-episode results for all four methods
    and a 6-condition DDIM step-count ablation over 30 episodes each.
  \item \textbf{Flow Matching superiority:} FM achieves 98\% success at 8.9\,ms/step,
    outperforming DDIM in both accuracy and speed, and beating DDPM by 18 pp.
  \item \textbf{5-step efficiency:} DDIM with 5 steps achieves 100\% success,
    demonstrating practical viability for real-time robot control.
\end{itemize}

% ============================================================
\section{Related Work}
% ============================================================

\subsection{Behavioral Cloning and Its Limitations}

Behavioral cloning treats demonstration data as supervised (observation, action)
pairs and trains a regressor~\citep{pomerleau1989alvinn}.
While effective when action distributions are unimodal, BC suffers from two key
failures: (a)~mode averaging in multi-modal settings and (b)~compounding errors as
the agent moves into states not covered by training.
\citet{florence2022ibc} showed that implicit behavior cloning (IBC) addresses mode
averaging through energy-based models, but at high inference cost.
Our BC baseline isolates the mode-averaging problem directly.

\subsection{Denoising Diffusion Models}

\citet{ho2020ddpm} introduced DDPM, defining a Markov forward process that
gradually corrupts data with Gaussian noise and trains a neural network to reverse
it.
The training objective is noise prediction:
$\mathcal{L} = \mathbb{E}\bigl[\|\boldsymbol{\varepsilon}_\theta(\mathbf{x}_t, t) -
\boldsymbol{\varepsilon}\|^2\bigr]$.
\citet{song2021ddim} subsequently showed that DDPM training implicitly learns a
score function~$\nabla_{\mathbf{x}} \log p(\mathbf{x}_t)$, enabling a deterministic
ODE (DDIM) that samples in far fewer steps.
Since its introduction, diffusion has driven state-of-the-art in image generation
(DALL$\cdot$E 2, Stable Diffusion), audio, video, and---as in our work---robot
control.

\subsection{Flow Matching}

\citet{lipman2023flow} and \citet{liu2023rectified} independently proposed Flow
Matching (FM) as a cleaner alternative.
Instead of learning to denoise a stochastically corrupted signal, FM learns to
follow a straight-line flow from noise to data:
$\mathbf{x}_t = (1{-}t)\,\mathbf{x}_0 + t\,\boldsymbol{\varepsilon}$.
The target velocity is $\mathbf{u} = \boldsymbol{\varepsilon} - \mathbf{x}_0$,
making training strictly supervised at every step.
FM has demonstrated superior sample quality at equivalent inference cost, with
naturally well-conditioned ODE trajectories.

\subsection{Diffusion Policy}

\citet{chi2023diffusion} applied diffusion to robot control by framing action
generation as a denoising problem conditioned on observations.
Their key insight is that FiLM-conditioned U-Nets are ideally suited to this
conditional generation task.
They demonstrated that diffusion policies outperform BC, LSTM-GMM, and IBC on
multiple manipulation tasks.
Our work reproduces and extends their PushT results with an additional FM baseline
and a systematic inference-step study.

% ============================================================
\section{Task and Dataset}
% ============================================================

\subsection{The PushT Task}

PushT is a 2D simulated robot manipulation task.
A circular end-effector (the agent) moves in a bounded 2D plane;
a T-shaped block sits on a table; and a gray \textit{target zone} marks the
desired block pose.
An episode succeeds when block-target overlap exceeds 90\%, and terminates at
success or 300 steps.

The task is deceptively difficult: the agent must contact the correct face of the
T-block (not the stem), approach from the correct angle, and reposition when the
block overshoots.
Crucially, the optimal approach direction is \emph{bimodal}---push from the left or
right are equally valid---making PushT an ideal testbed for multi-modal policies.

\subsection{Dataset}

We use the Columbia PushT expert demonstration dataset~\citep{chi2023diffusion},
the same dataset used in the original paper.
Key statistics are summarized in Table~\ref{tab:dataset}.

\begin{table}[h]
  \centering
  \caption{PushT expert demonstration dataset statistics.}
  \label{tab:dataset}
  \small
  \begin{tabular}{@{}ll@{}}
    \toprule
    \textbf{Property} & \textbf{Value} \\
    \midrule
    Format            & Zarr archive (\texttt{pusht\_cchi\_v7\_replay.zarr}) \\
    Episodes          & 206 \\
    Total timesteps   & 25,650 \\
    Avg.\ episode length & $\sim$124 steps \\
    Observation dims  & 5 (agent $x,y$; block $x,y,\theta$) \\
    Action dims       & 2 (velocity $v_x, v_y$) \\
    Coordinate range  & 0--512 pixels \\
    \bottomrule
  \end{tabular}
\end{table}

\subsection{Data Preprocessing}

\paragraph{Normalization.}
All observations and actions are scaled to $[-1, 1]$ using per-dimension
min-max normalization.
This is essential: the DDPM noise schedule assumes unit-scale data, and
pixel-coordinate values (0--512) would render the noise magnitude
completely wrong.

\paragraph{Sliding-window sampling.}
We train on overlapping short sequences rather than single~(obs,\;action) pairs.
The model observes the last $\Tobs = 2$ timesteps and predicts the next $\Tpred = 16$
actions; only the first $\Tact = 8$ are executed before replanning.
Episodes are padded at boundaries so every timestep is a valid window center,
producing approximately \textbf{25,000 training samples} from 206 episodes.

\begin{figure}[t]
  \centering
  \includegraphics[width=\linewidth]{multimodal_motivation.png}
  \caption{%
    \textbf{Multi-modal motivation.}
    \emph{Left:} Expert demonstration trajectories on PushT showing two distinct
    push strategies (left approach vs.\ right approach).
    \emph{Right:} Bimodal action distribution from the same observation state.
    A BC regressor averages these two modes, producing an action ($\times$) that
    commits to neither strategy and fails the task.
  }
  \label{fig:multimodal}
\end{figure}

Figure~\ref{fig:multimodal} illustrates why BC fails on this task.
The bimodal action distribution is the core empirical motivation for using a
distributional policy.

% ============================================================
\section{Methods}
% ============================================================

\subsection{Problem Formulation}

At each control step, the agent observes a history of $\Tobs = 2$ normalized state
vectors (10 values total).
The policy must produce an action chunk of $\Tpred = 16$ actions.
We model this as learning the conditional distribution
$p(\mathbf{a}_{0:\Tpred} \mid \mathbf{o}_{-\Tobs+1:0})$.
Diffusion models approximate this via iterative denoising, allowing samples that
faithfully reflect all modes of the expert distribution rather than collapsing to
the mean.

\subsection{Network Architecture: ConditionalUnet1D}
\label{sec:arch}

The backbone is a \textbf{1D temporal U-Net} (68.95M parameters) that treats the
action sequence $(\Tpred \times d_{\text{act}}) = (16 \times 2)$ as a 1D signal
with 2 channels.
The encoder--decoder structure with skip connections allows the model to mix
fine-grained action values with coarse temporal structure.

\paragraph{Conditioning pipeline.}
Two parallel encoders produce the 512-dim conditioning vector:
\begin{itemize}
  \item \emph{Observation encoder:} A fully learnable 2-layer MLP maps the flattened
    observation history (10 values) to 256 dimensions:
    $\text{Linear}(10 \!\to\! 256) \to \text{Mish} \to \text{Linear}(256 \!\to\! 256)$.
    All 68,864 parameters are trained end-to-end.
  \item \emph{Timestep encoder:} A hybrid pipeline combines a fixed sinusoidal
    positional embedding (256-dim, no gradients) with a learnable MLP:
    $\text{Linear}(256 \!\to\! 1024) \to \text{Mish} \to \text{Linear}(1024 \!\to\! 256)$.
  \item \emph{Conditioning:} The two 256-dim embeddings are concatenated into a
    512-dim vector passed to every residual block.
\end{itemize}

\paragraph{U-Net channel progression.}
\vspace{2pt}
\begin{center}
\small
\begin{tabular}{@{}llll@{}}
  \toprule
  \textbf{Stage} & \textbf{Input} & \textbf{Output} & \textbf{Operation} \\
  \midrule
  Down 1 & $(2, 16)$ & $(256, 16)$ & ResBlock $\times$ 2 \\
  Down 2 & $(256, 16)$ & $(256, 8)$  & Downsample \\
  Down 3 & $(256, 8)$  & $(512, 8)$  & ResBlock $\times$ 2 \\
  Down 4 & $(512, 8)$  & $(512, 4)$  & Downsample \\
  Bottleneck & $(512, 4)$ & $(1024, 4)$ & ResBlock $\times$ 2 \\
  Up 1   & $(1024{+}1024, 4)$ & $(512, 4)$ & ResBlock $\times$ 2 \\
  Up 2   & $(512, 4)$  & $(512, 8)$  & Upsample \\
  Up 3   & $(512{+}256, 8)$  & $(256, 8)$ & ResBlock $\times$ 2 \\
  Up 4   & $(256, 8)$  & $(256, 16)$ & Upsample \\
  Head   & $(256, 16)$ & $(2, 16)$   & Conv1d $1{\times}1$ \\
  \bottomrule
\end{tabular}
\end{center}
\vspace{4pt}

Each \textbf{ConditionalResidualBlock1D} follows the pattern:
Conv1d $\to$ GroupNorm $\to$ FiLM $\to$ Mish $\to$ Conv1d $\to$ GroupNorm $\to$ FiLM $\to$ Mish,
with a $1{\times}1$ residual Conv1d handling channel mismatches.
All convolutions use kernel size 5 with padding 2 to preserve sequence length.

\subsection{FiLM Conditioning}

Feature-wise Linear Modulation~\citep{perez2018film} injects the conditioning
signal into each feature map:
\begin{equation}
  \text{FiLM}(\mathbf{x},\, \mathbf{c})
  = \gamma(\mathbf{c}) \odot \text{GroupNorm}(\mathbf{x})
  \;+\; \beta(\mathbf{c}),
  \label{eq:film}
\end{equation}
where $\gamma$ and $\beta$ are computed via
$\text{Linear}(512, 2 \cdot d_{\text{out}})$, then split.
FiLM is strictly more expressive than input concatenation because it modulates
every feature map \emph{independently} based on the current observation and noise
level.

\subsection{DDPM: Stochastic Denoising}

DDPM~\citep{ho2020ddpm} defines a forward process that adds Gaussian noise over
$K = 100$ discrete steps using a \textit{cosine schedule}:
\begin{align}
  \mathbf{a}_k &= \sqrt{\alphabar_k}\,\mathbf{a}_0 \;+\; \sqrt{1 - \alphabar_k}\,\boldsymbol{\varepsilon},
  \quad \boldsymbol{\varepsilon} \sim \mathcal{N}(\mathbf{0}, \mathbf{I}), \label{eq:fwd} \\
  \alphabar_k &= \cos^2\!\Bigl(\tfrac{k/K + 0.008}{1.008}\cdot\tfrac{\pi}{2}\Bigr). \label{eq:cosine}
\end{align}
At $k{=}0$: $\alphabar_0 \approx 1.0$ (clean). At $k{=}99$: $\alphabar_{99} \approx 2{\times}10^{-8}$
(near-pure noise).
The model is trained to predict the added noise:
\begin{equation}
  \mathcal{L}_{\text{DDPM}} = \mathbb{E}\bigl[
    \|\boldsymbol{\varepsilon}_\theta(\mathbf{a}_k, k, \mathbf{o}) - \boldsymbol{\varepsilon}\|^2
  \bigr].
  \label{eq:ddpm_loss}
\end{equation}
Inference runs 100 reverse steps from $\mathbf{a}_{99} \sim \mathcal{N}(\mathbf{0}, \mathbf{I})$:
\begin{equation}
  \mathbf{a}_{k-1} = \frac{1}{\sqrt{\alpha_k}}\!\left(\mathbf{a}_k - \frac{\beta_k}{\sqrt{1-\alphabar_k}}
  \,\boldsymbol{\varepsilon}_\theta\right) + \sqrt{\beta_k}\,\mathbf{z},
  \quad \mathbf{z} \sim \mathcal{N}(\mathbf{0},\mathbf{I}) \text{ if } k{>}0.
  \label{eq:ddpm_rev}
\end{equation}

\subsection{DDIM: Deterministic Fast Sampling}

DDIM~\citep{song2021ddim} reuses the trained DDPM model but replaces the stochastic
reverse process with a deterministic ODE, enabling inference in as few as 10 steps.
Setting $\eta = 0$ (no stochastic noise), the update at each step is:
\begin{align}
  \hat{\mathbf{a}}_0 &= \frac{\mathbf{a}_t - \sqrt{1{-}\alphabar_t}\,\boldsymbol{\varepsilon}_\theta}
    {\max\!\bigl(\sqrt{\alphabar_t},\; 10^{-3}\bigr)}, \label{eq:ddim1} \\
  \hat{\mathbf{a}}_0 &\leftarrow \operatorname{clip}(\hat{\mathbf{a}}_0,\, -1,\, 1), \label{eq:ddim_clip} \\
  \mathbf{a}_{t_{\text{prev}}} &= \sqrt{\alphabar_{t_{\text{prev}}}}\,\hat{\mathbf{a}}_0
    + \sqrt{1{-}\alphabar_{t_{\text{prev}}}}\,\boldsymbol{\varepsilon}_\theta. \label{eq:ddim2}
\end{align}
The clamping in Eq.~\eqref{eq:ddim1} is \emph{critical}: without it,
$\sqrt{\alphabar_{99}} \approx 1.6{\times}10^{-4}$ causes a division by near-zero
that blows up $\hat{\mathbf{a}}_0$ to ${\pm}10^6$, producing a 0\% success rate despite a
well-trained model. We detail this in Section~\ref{sec:challenges}.

\subsection{Flow Matching}

Flow Matching~\citep{lipman2023flow,liu2023rectified} replaces the stochastic
diffusion process with a straight-line interpolation between data and noise:
\begin{equation}
  \mathbf{x}_t = (1{-}t)\,\mathbf{a}_0 \;+\; t\,\boldsymbol{\varepsilon},
  \quad t \in [0, 1].
  \label{eq:fm_interp}
\end{equation}
The model predicts the \textit{velocity} (direction of flow in action space):
\begin{equation}
  \mathbf{u} = \boldsymbol{\varepsilon} - \mathbf{a}_0, \qquad
  \mathcal{L}_{\text{FM}} = \mathbb{E}\bigl[\|\mathbf{v}_\theta(\mathbf{x}_t, t, \mathbf{o}) - \mathbf{u}\|^2\bigr].
  \label{eq:fm_loss}
\end{equation}
Inference runs the Euler ODE backward from $t{=}1$ (pure noise) to $t{=}0$
(clean action) in $N$ steps:
\begin{equation}
  \mathbf{x}_{t - \Delta t} = \mathbf{x}_t - \Delta t \cdot \mathbf{v}_\theta(\mathbf{x}_t, t, \mathbf{o}),
  \quad \Delta t = 1/N.
  \label{eq:fm_euler}
\end{equation}
FM uses the \emph{same U-Net architecture} as DDPM; only the training loss and
inference procedure change.
The continuous time $t \in [0,1]$ is scaled by 100 before passing to the
sinusoidal timestep embedding.
Crucially, FM requires \emph{no noise schedule tuning}---the straight-line
interpolation has no free parameters.

\subsection{Training Procedure}

\paragraph{Training loop.}
For each batch: (1)~sample noise $\boldsymbol{\varepsilon} \sim \mathcal{N}(\mathbf{0}, \mathbf{I})$ and
timestep $k \sim \mathcal{U}\{0, \ldots, K{-}1\}$; (2)~compute noisy action via
Eq.~\eqref{eq:fwd} (DDPM) or Eq.~\eqref{eq:fm_interp} (FM); (3)~forward pass;
(4)~compute MSE loss; (5)~backpropagate, clip gradients at max norm 1.0;
(6)~AdamW step; (7)~EMA update: $\theta_{\text{ema}} \leftarrow 0.995\,\theta_{\text{ema}} + 0.005\,\theta_{\text{train}}$.

\paragraph{Hyperparameters.}
All models train for 300 epochs on an NVIDIA A100 GPU via the Northeastern Explorer
HPC cluster (SLURM scheduler).
Key hyperparameters are summarized in Table~\ref{tab:hparams}.

\begin{table}[h]
  \centering
  \caption{Training hyperparameters (shared across all diffusion methods).}
  \label{tab:hparams}
  \small
  \begin{tabular}{@{}ll@{}}
    \toprule
    \textbf{Hyperparameter} & \textbf{Value} \\
    \midrule
    Architecture            & ConditionalUnet1D (FiLM) \\
    Parameters              & 68.95M \\
    Channel dims (down)     & (256, 512, 1024) \\
    Conditioning dim        & 512 (256 obs + 256 timestep) \\
    $\Tobs / \Tpred / \Tact$ & 2 / 16 / 8 \\
    Diffusion steps $K$     & 100 (DDPM/DDIM); continuous (FM) \\
    Noise schedule          & Cosine (DDPM/DDIM); linear interp.\ (FM) \\
    Epochs                  & 300 \\
    Batch size              & 256 \\
    Optimizer               & AdamW \\
    Peak learning rate      & $1{\times}10^{-4}$ \\
    LR schedule             & Cosine decay + 500-step warmup \\
    Gradient clipping       & max norm = 1.0 \\
    EMA decay               & 0.995 \\
    Hardware                & NVIDIA A100 / V100 SXM2 \\
    \bottomrule
  \end{tabular}
\end{table}

\paragraph{Learning rate schedule.}
The LR follows a two-phase schedule:
(1)~\emph{linear warmup} from 0 to $10^{-4}$ over 500 steps, preventing large
gradient updates during early random-weight training;
(2)~\emph{cosine decay} from $10^{-4}$ to $\approx$0 over the remaining steps.

\subsection{Receding-Horizon Control}

At evaluation, the agent replans every $\Tact = 8$ steps:
\begin{enumerate}[topsep=3pt]
  \item Observe the last $\Tobs = 2$ states from a sliding deque.
  \item Normalize observations; run the full denoising process to produce
    $\Tpred = 16$ normalized actions.
  \item Unnormalize to pixel-space velocities; execute the first $\Tact = 8$
    actions in the environment.
  \item Append new observations to the deque (oldest drops off). Repeat.
\end{enumerate}
Predicting a longer horizon ($\Tpred = 16$) than is executed ($\Tact = 8$)
encourages temporal consistency---a form of implicit model predictive control.

% ============================================================
\section{Experiments}
% ============================================================

\subsection{Evaluation Setup}

All methods are evaluated in a \texttt{gymnasium}-compatible PushT simulation
(\texttt{gym\_pusht}).
Each episode begins with a randomly positioned T-block and fixed agent start.
We report: \emph{success rate} (coverage ${\ge}0.9$), \emph{mean coverage},
\emph{mean episode length}, and \emph{inference time per control step}.

\paragraph{BC baseline.}
A 2-hidden-layer MLP (135K parameters, hidden dim 256, Mish activations) trained
for 200 epochs.
This policy performs standard regression from the flattened observation history to
a single action, serving as the mode-averaging failure-mode reference.

\begin{figure}[t]
  \centering
  \includegraphics[width=0.88\linewidth]{diffusion_forward_process.png}
  \caption{%
    \textbf{DDPM forward process.}
    Clean action sequences ($t{=}0$, left) are progressively corrupted with
    Gaussian noise over 100 timesteps until reaching pure noise ($t{=}99$, right).
    Training reverses this process: the U-Net learns to recover coherent action
    sequences from noise, conditioned on the current observation.
  }
  \label{fig:forward}
\end{figure}

\subsection{Training Dynamics}

Figure~\ref{fig:training} shows loss curves for DDPM and Flow Matching over 300
epochs.
DDPM achieves a final noise-prediction MSE of \textbf{0.0055}; FM achieves a final
velocity-prediction MSE of \textbf{0.0141}.
These values are \emph{not directly comparable} (noise MSE vs.\ velocity MSE),
but both exhibit a sharp drop in the first 10 epochs followed by smooth convergence.
FM's higher absolute loss does not indicate worse learning---FM achieves 98\%
success vs.\ DDPM's 80\%.

\begin{figure}[t]
  \centering
  \includegraphics[width=\linewidth]{ddpm_vs_fm_300ep.png}
  \caption{%
    \textbf{Training loss curves.}
    DDPM (blue) and Flow Matching (orange) over 300 epochs.
    Both converge smoothly. The losses are not directly comparable (noise MSE vs.\
    velocity MSE), but both reach low values with a cosine LR schedule.
  }
  \label{fig:training}
\end{figure}

\subsection{Main Results}

Table~\ref{tab:main} presents the main evaluation results over 50 episodes per
method.
Figure~\ref{fig:results} visualises the success rates.

\begin{table}[h]
  \centering
  \caption{Main evaluation results (50 episodes per method). Paper baselines are
    approximate read-offs from~\citet{chi2023diffusion} figures.
    $\star$ = best result.}
  \label{tab:main}
  \small
  \setlength{\tabcolsep}{5pt}
  \begin{tabular}{@{}lccccr@{}}
    \toprule
    \textbf{Method} & \textbf{Eps.} & \textbf{Success} & \textbf{Cov.} & \textbf{ms/step} & \textbf{Steps} \\
    \midrule
    \rowcolor{BestGreen}
    \textbf{Flow Matching (ours)} & 50 & \textbf{98\%}~$\star$ & 0.989 & 8.9  & 10 \\
    \textbf{DDIM (ours, 300ep)}   & 50 & \textbf{92\%}        & 0.969 & 10.0 & 10 \\
    \rowcolor{RowAlt}
    \textbf{DDPM (ours, 300ep)}   & 50 & \textbf{80\%}        & 0.830 & 84.0 & 100 \\
    \textbf{BC baseline (ours)}   & 50 & \textbf{4\%}         & 0.241 & 0.6  & --- \\
    \midrule
    \rowcolor{RowAlt}
    DDIM \citep{chi2023diffusion} & --- & $\sim$90\%          & ---   & ---  & 10 \\
    DDPM \citep{chi2023diffusion} & --- & $\sim$92\%          & ---   & ---  & 100 \\
    \bottomrule
  \end{tabular}
\end{table}

\begin{figure}[h]
  \centering
  \includegraphics[width=0.95\linewidth]{final_comparison.png}
  \caption{%
    \textbf{Main results.}
    Success rates across all four methods (50 episodes each).
    All three diffusion variants dramatically outperform BC (4\%), with FM
    achieving 98\%.
    The 24$\times$ gap between BC and FM validates the core motivation for
    distributional action prediction.
  }
  \label{fig:results}
\end{figure}

\paragraph{Flow Matching vs.\ DDIM.}
FM outperforms DDIM by 6 pp (98\% vs.\ 92\%) at comparable speed (8.9\,ms vs.\
10\,ms).
We attribute this to FM's \emph{straight-line ODE path}: the model predicts a
constant velocity directly connecting noise to data, avoiding the curved cosine
trajectory that allocates disproportionate network capacity to the high-noise
regime.

\paragraph{DDPM vs.\ DDIM.}
DDIM outperforms DDPM (92\% vs.\ 80\%) while being 10$\times$ faster
(10\,ms vs.\ 84\,ms).
The deterministic reverse process eliminates stochastic noise injected at each of
DDPM's 100 steps---noise that occasionally pushes trajectories away from
high-probability action modes.

\paragraph{BC baseline: the 24$\times$ gap.}
The BC baseline achieves only 4\% success.
Every episode hits the 300-step timeout (mean length = 300.0, mean coverage = 0.241).
The MLP converges to the mode-averaged action, oscillating in place without
committing to either push strategy.
This \textbf{24$\times$ performance gap} (4\% vs.\ 98\%) is the core empirical
result confirming that multi-modal expert demonstrations require a distributional
policy.

\paragraph{Comparison to~\citeauthor{chi2023diffusion}.}
Our DDIM result (92\%) matches the paper's DDIM ($\sim$90\%) within the 95\% CI of
$\pm$5.5 pp for 50 episodes.
Our FM result (98\%) establishes a new reference---Chi et~al.\ did not report FM on
PushT.
Our DDPM result (80\% vs.\ $\sim$92\%) is below the paper, likely due to
differences in random seed variance over 50 episodes and potential training
duration differences.

\subsection{DDIM Inference Steps Ablation}

DDIM's key advantage is inference-time flexibility: the same trained model can use
any step count without retraining.
Table~\ref{tab:ablation} and Figure~\ref{fig:ablation} sweep $N \in \{1, 5, 10, 20, 50, 100\}$
over 30 episodes each using the same 300-epoch DDPM checkpoint.

\begin{table}[h]
  \centering
  \caption{DDIM inference steps ablation (30 episodes per condition).
    $\star$ = best efficiency point.}
  \label{tab:ablation}
  \small
  \setlength{\tabcolsep}{4pt}
  \begin{tabular}{@{}rcccl@{}}
    \toprule
    \textbf{Steps} & \textbf{Success} & \textbf{Coverage} & \textbf{ms/step} & \textbf{Notes} \\
    \midrule
    1             & 0\%              & 0.088             & 2.3              & Total failure \\
    \rowcolor{BestGreen}
    5             & \textbf{100\%}~$\star$ & 0.991      & 5.6              & Best efficiency \\
    10            & 93.3\%           & 0.958             & 9.8              & $\leftarrow$ Our default \\
    \rowcolor{RowAlt}
    20            & 100\%            & 0.996             & 18.1             & Marginally better \\
    50            & 93.3\%           & 0.947             & 42.9             & No gain over 10 \\
    \rowcolor{RowAlt}
    100           & 100\%            & 0.989             & 84.0             & Full DDPM cost \\
    \bottomrule
  \end{tabular}
\end{table}

\begin{figure}[h]
  \centering
  \includegraphics[width=\linewidth]{steps_ablation.png}
  \caption{%
    \textbf{DDIM steps ablation.}
    Success rate (bars, left axis) and inference time per control step (line,
    right axis) as a function of denoising steps.
    5 steps achieves 100\% success at 5.6\,ms---the optimal operating point.
    Diminishing returns set in beyond 10 steps.
  }
  \label{fig:ablation}
\end{figure}

\paragraph{Key findings.}
\textbf{(1)~1 step fails completely (0\%):} A single denoising step cannot
reconstruct a coherent action sequence from pure noise---iterative refinement is
genuinely necessary.
\textbf{(2)~5 steps achieves 100\%:} The model's learned score function is
well-conditioned enough that 5 refinement steps fully recover high-quality actions,
demonstrating practical viability at 5.6\,ms/step.
\textbf{(3)~Diminishing returns beyond 10 steps:} Going from 10 to 100 steps adds
74\,ms of latency without a systematic success-rate improvement.
The sweet spot for real-robot deployment is \textbf{5--10 steps}.

% ============================================================
\section{Key Implementation Challenges}
\label{sec:challenges}
% ============================================================

\subsection{The DDIM Numerical Stability Fix}

\paragraph{Problem.}
The cosine schedule at $K{=}100$ yields $\alphabar_{99} \approx 2{\times}10^{-8}$.
DDIM's first inference step divides by
$\sqrt{\alphabar_{99}} \approx 1.6{\times}10^{-4}$,
producing predicted clean actions on the order of $\pm 10^6$.
Since normalized actions lie in $[-1, 1]$, these values trigger environment
velocity clipping to zero---the agent freezes.

\paragraph{Symptom.}
Despite a well-trained model (training loss = 0.013), the evaluated success rate
was \textbf{0.000}.
Every GIF frame showed an identical frozen agent---visually indistinguishable from
an untrained model.

\paragraph{Fix.}
\begin{lstlisting}[language=Python, caption={DDIM numerical stability fix (two lines).}]
# Before: numerical explosion with cosine schedule at K=100
a0_pred = (noisy_a - (1 - ab_t).sqrt() * eps_pred) / ab_t.sqrt()

# After: stable prediction with clamped denominator + clipping
a0_pred = (noisy_a - (1 - ab_t).sqrt() * eps_pred) \
            / ab_t.sqrt().clamp(min=1e-3)
a0_pred = a0_pred.clamp(-1.0, 1.0)
\end{lstlisting}

\paragraph{Outcome.}
Success rate jumped from \textbf{0\% to 92\%} on the same checkpoint, matching the
HuggingFace Diffusers reference and the original Chi et al.\ codebase.
The lesson: numerical stability issues in diffusion samplers can cause complete
\emph{silent} failure---the model trains perfectly but produces garbage at
inference time.

\subsection{Zarr API Breaking Change}

The PushT dataset is stored in Zarr format.
Between Zarr v2 and v3, the positional \texttt{open(path, mode)} signature was
changed to require keyword arguments (\texttt{open(store=path, mode='r')}).
The v3 API silently accepted the positional call but returned an empty store,
causing training to start normally but on zero data.
We identified this through unit tests on the dataset class and pinned
\texttt{zarr>=2.14,<4} in \texttt{requirements.txt}.

\subsection{Training Crash and Checkpoint Recovery}

The initial training run crashed at epoch 28 due to machine hibernation, with no
checkpoint saved (default interval: 50 epochs).
We added \texttt{--save\_interval} as a CLI argument (default: 10 epochs) and
restarted.
This experience informed our final strategy: checkpoint every 10 epochs plus a
\texttt{best.pt} overwritten whenever evaluation loss improves.

% ============================================================
\section{Discussion}
% ============================================================

\paragraph{Why does Flow Matching win?}
FM's straight-line interpolation
$\mathbf{x}_t = (1{-}t)\mathbf{a}_0 + t\boldsymbol{\varepsilon}$
creates a fundamentally different geometry than DDPM's cosine schedule.
In DDPM, most SNR change occurs at small~$t$ (high-SNR regime), creating a curved
trajectory that the network must learn to navigate.
FM's constant-velocity field is always aimed directly from noise to data---easier
to learn and easier for the ODE solver to integrate.
This manifests as both higher success rates (98\% vs.\ 92\%) and fewer required
steps.

\paragraph{Temporal coherence vs.\ multi-modality.}
The 5-step DDIM result (100\% success via deterministic sampling) suggests the
model is not relying on fine-grained stochastic exploration.
This implicates \emph{temporal coherence} as the primary mechanism: predicting 16
actions jointly guarantees that consecutive actions are internally consistent.
BC predicts one action at a time, producing independent predictions that may jump
between modes mid-episode.

\paragraph{Practical deployment implications.}
\begin{itemize}
  \item \textbf{FM or DDIM, not DDPM:} DDPM's 84\,ms/step latency is a practical
    bottleneck. FM at 8.9\,ms and DDIM at 10\,ms are viable in a 30\,Hz control loop.
  \item \textbf{5 steps may suffice:} If latency is critical (e.g., reactive safety
    behaviors), 5 DDIM steps at 5.6\,ms may be the optimal operating point.
  \item \textbf{EMA is essential:} Evaluating with EMA weights (decay = 0.995,
    effective window = 200 gradient steps) consistently outperformed non-EMA
    evaluation by filtering out noisy weight updates without sacrificing
    final performance.
\end{itemize}

% ============================================================
\section{Conclusion}
% ============================================================

We presented a complete from-scratch implementation of Diffusion Policy for robot
manipulation on PushT.
Our evaluation of four methods confirms that distributional action policies
dramatically outperform direct regression in multi-modal expert settings.
The \textbf{24$\times$ performance gap} between BC (4\%) and Flow Matching (98\%)
is a direct empirical demonstration of the mode-averaging failure of BC.

Among diffusion methods, Flow Matching is the best choice: no noise schedule
tuning, a strictly supervised velocity-matching objective, and the highest success
rate (98\%) at the lowest inference latency (8.9\,ms).

Our most important engineering contribution is the DDIM numerical stability fix
(clamping the cosine-schedule denominator at $10^{-3}$ and clipping predictions
to $[-1, 1]$), which lifts DDIM from 0\% to 92\% on an otherwise well-trained
model---a critical finding for practitioners implementing diffusion policies.

Finally, our inference-step ablation reveals that 5 denoising steps achieve 100\%
success---opening the door to real-time deployment at 5\,ms per control step,
well within the latency budget of physical robot systems.

% ── ACKNOWLEDGEMENTS ─────────────────────────────────────────
\paragraph{Acknowledgements.}
We thank the Northeastern University Research Computing team for access to the
Explorer HPC cluster (NVIDIA A100 and V100 SXM2 nodes) used to train all models.
We also thank \citet{chi2023diffusion} for open-sourcing the Columbia PushT dataset
and reference codebase.

% ── REFERENCES ───────────────────────────────────────────────
\bibliographystyle{plainnat}
\bibliography{refs}

\end{document}
